\section{Conclusion}
\label{sec:conclusion}
This paper presented \textbf{MPTA-Repair}, an automatic repair workflow
for multi-property, multi-counterexample clock-bound repair of
timed-automata control models. The method addresses a practical
limitation of local TDT-based repair: control models are commonly
validated against multiple timing properties, and one timing defect may
produce several related counterexamples. Repairing one trace in
isolation may therefore be insufficient.

MPTA-Repair collects diagnostic traces from violated properties, uses
lightweight relations among properties, traces, and repairable bounds to
guide repair search, and synthesizes small clock-bound edits using a
MaxSMT-based backend. Every candidate is validated by full-property
regression and a TARTAR-style admissibility check based on functional
equivalence \cite{koelbl2019clock}. This design ensures that accepted repairs are not
merely local counterexample fixes, but system-level timing corrections
that preserve intended behavior.

Our evaluation on non-control and control-domain timed-automata
benchmarks shows that MPTA-Repair improves repair success over an
iterative TARTAR-style baseline, especially on control-domain models
where multi-property interaction and admissibility constraints are
prominent. The results also show that failures remain meaningful:
unrepaired cases often require changes outside the clock-bound repair
space or exceed the configured solving and validation budgets.

Future work will extend the repair space beyond numerical clock bounds
to include resets, clock references, synchronization labels, and
structural edits. We also plan to strengthen counterexample generation,
improve relation-guided decomposition, and evaluate the method on larger
industrial timed-automata models.
